В проекте используются два типа скоров: docking-скоры $d_A, d_B$ (UniDock GPU, kcal/mol)
служат \textbf{RL reward} для генерации, а Boltz-2 affinity $a_A, a_B$ ($\log_{10}$\,(IC50-like), $\mu$M;
обучена на смеси Ki/Kd/IC50/EC50)
--- \textbf{ортогональная evaluation surrogate}. Два скорера не смешиваются;
$a_X$ интерпретируется как ранжирующая метрика, а не калиброванный IC50.

Генерация выполняется в трёх режимах:
\texttt{single\_a} (reward $= -d_A$), \texttt{single\_b} (reward $= -d_B$),
\texttt{double} (reward $= -\max(d_A, d_B)$).
Reward backend --- UniDock GPU (\texttt{unidock\_gpu}), GPU-ускоренный докинг,
CLI-совместимый с AutoDock Vina.

Скоринг выполняется скриптом \texttt{scripts/score.py}: Boltz-2 оценивает каждую
сгенерированную молекулу по двум мишеням пары с записью \texttt{scored\_joined.csv}.
Далее метрики (hit-rate по $a_A, a_B$) и статистика рассчитываются в
\texttt{scripts/compute\_metrics.py} с двухуровневым bootstrap и поправкой на множественные сравнения.

Оркестрация разделена на два этапа:
\begin{enumerate}
\item Stage 1: \texttt{scripts/run\_full\_experiment.py} (generate + score).
\item Stage 2: \texttt{scripts/run\_post\_pipeline.py} (aggregate, metrics, plots, render, summary, LaTeX).
\end{enumerate}

\subsection*{Reward function для double-режима}

Ключевой design choice --- как FREED++ комбинирует $d_A$ и $d_B$ в один reward.

\textbf{Primary experiment:} $R = -\max(d_A, d_B)$ (worst-case).
Молекула оценивается по худшему docking-скору; гарантирует активность к обоим таргетам.

\textbf{Sensitivity analysis (offline):} $R = -(d_A + d_B)$ (сумма) и другие
скаляризации (geometric, harmonic, Chebyshev) пересчитываются по уже сгенерированным
молекулам. Ближайший прецедент суммы: SFG-Drug (MDPI 2025) --- dual MEK1/mTOR.

\subsection*{Категории пар}

Эксперимент включает 9 пар мишеней (8 активных, 1 заблокированная),
классифицированных по трём осям: \emph{disease context} (AD / Non-AD),
\emph{structural relation} (Homo / Hetero) и \emph{cofactor presence} (Clean / Present).
Пары разделены на заранее зафиксированные множества:

\begin{itemize}
\item \textbf{Core set} (5 пар) --- основная проверка гипотезы:
  CDK2\_ERK2 (Non-AD, Homo), AChE\_BChE (AD, Homo),
  BACE1\_GSK3B, AChE\_BACE1, GSK3B\_CDK5 (AD, Hetero).
  Все без кофакторов; Boltz-2 affinity валиден.
\item \textbf{Auxiliary set} (3 пары) --- robustness и sensitivity:
  EGFR\_HER2 (второй Homo-контроль), BACE1\_CDK5 (замена tau-киназы),
  AChE\_MAOB (cofactor-sensitive: FAD removed in receptor prep, Boltz-2 binary only).
\item \textbf{Blocked}: COX1\_COX2 (non-human PDB: COX1 Ovis aries, COX2 Mus musculus).
\end{itemize}

Гомологичные пары (AChE\_BChE, CDK2\_ERK2) служат ``лёгким'' тестом --- высокое сходство
карманов облегчает dual generation. Гетерогенные AD-пары (BACE1\_GSK3B, GSK3B\_CDK5) ---
``жёсткий'' тест, проверяющий способность \texttt{double} режима находить dual-active
молекулы для структурно различных мишеней.

Деление core/auxiliary зафиксировано до начала экспериментов; primary statistical
tests проводятся только на core set (подробнее: \texttt{docs/selection\_criteria.md}).
